\section{Introduction}

Low-light image enhancement is the task of transforming an under-exposed image into a more visible and visually useful one. This problem is important in several practical scenarios, such as night photography, surveillance, robotics, and any visual system that must operate under poor illumination. However, low-light enhancement is not simply a matter of increasing brightness. If the enhancement is too aggressive, the output can become over-exposed, noisy, desaturated, or affected by artificial halos. If it is too conservative, important regions can remain dark. A good model should therefore recover visibility while preserving structure, colors, edges, and realistic image content.

A central difficulty of this task is generalization. A model trained on one dataset may perform well on images captured under similar conditions, but fail when the camera response, illumination distribution, noise level, or scene content changes. This is especially relevant for low-light enhancement because different datasets may represent darkness in very different ways. For this reason, evaluating only on images from the same dataset and acquisition setting used for training is not sufficient to understand whether a model has learned a robust enhancement strategy.

This work studies low-light image enhancement under a cross-dataset generalization protocol, using the datasets described in Section~\ref{sec:datasets}. All models are trained on the LOL-v2 Real dataset, which is used as the source domain, and are then evaluated on several test datasets without fine-tuning. The evaluation protocol includes the held-out LOL-v2 Real test set, the paired LOL-v2 Synthetic test set, MILLs, and the unpaired real-world stress-test dataset ExDark. LOL-v1 is also evaluated, but it is intentionally treated as a non-independent overlap-control set rather than as evidence of transfer: its inclusion is useful for documenting how the selected models behave when the same underlying low-light content is presented through the earlier benchmark organization. Together, these evaluation sets make it possible to distinguish in-domain behavior, transfer to synthetic degradation, robustness under controlled illumination levels, and behavior on real unpaired dark scenes.

The experimental protocol combines quantitative and qualitative evidence. Paired datasets are evaluated with full-reference metrics, while the unpaired ExDark dataset is evaluated with no-reference metrics and output diagnostics. In addition, all model families are inspected qualitatively through representative examples and failure cases. This combination is important because average scores alone do not fully explain how a model enhances low-light images: two outputs can have similar numerical quality while showing different behavior in terms of residual darkness, over-exposure, color cast, blur, noise, or halo artifacts.

The models are organized into four families. The first one is a classical U-Net, used as a strong and interpretable encoder-decoder reference. The second one is an illumination-guided U-Net, which adds an estimated illumination map as a fourth input channel in order to test whether explicit local brightness information improves enhancement. The third one is Frequency-Aware Exposure-Consistent Low-Light Enhancement (FAC-ECL), a training-objective variant that separates low-frequency illumination and color behavior from high-frequency detail behavior, while also encouraging stability under mild photometric perturbations. The fourth one is Exposure-Consistency-based Failure-Aware Refinement (ECL-FAR), an original refinement-based model family introduced in this project. ECL-FAR uses the model's own sensitivity to photometric perturbations as a cue for local refinement.

The goal is not only to rank these families, but also to understand what each design choice changes. In particular, this work investigates whether an explicit illumination cue is sufficient, whether frequency-aware supervision improves cross-dataset reconstruction, and whether photometric instability can be turned into useful information for refinement. The analysis shows that brightness control alone is not enough. The most effective families are those that also constrain how exposure, color, detail, and robustness interact.

The main contributions of this work are threefold. First, it implements a complete and reproducible low-light enhancement pipeline, including manifest generation, deterministic preprocessing, training, validation-based model selection, multi-seed evaluation, log aggregation, and qualitative export. Second, it defines and applies an evaluation protocol covering one held-out in-domain test set, three independent external datasets, and one deliberately non-independent LOL-v1 overlap-control set, using a total of 624 logged experimental runs to compare the selected model families under consistent conditions. Third, it develops the model comparison as an incremental design study: starting from a classical U-Net, testing whether explicit illumination guidance is sufficient, introducing frequency-aware and exposure-consistent supervision through FAC-ECL, and finally proposing ECL-FAR as an original instability-guided refinement family. Although ECL-FAR is only a first implementation and is not exhaustively optimized in all its components, it is already competitive with FAC-ECL, achieves the highest observed SSIM on MILLs, and is visually convincing in several of the final qualitative comparisons.

The full project is hosted in a GitHub repository, including source code, configuration files, reproducibility details, selected model configurations, and notes about dataset licenses and usage restrictions\footnote{\href{https://github.com/cosimo-curri/low-light-enhancement-cross-dataset}{Project Repository}}.

The rest of the report is organized as follows. The next section describes the datasets, licenses, and split construction. Then, the experimental protocol and the end-to-end pipeline are presented. The following sections analyze each model family separately, including architecture, objective, grid search, quantitative results, qualitative behavior, and failure mitigation strategies. After the family-specific analyses, an inter-family section compares the selected models directly and discusses the main trends observed across datasets. The report ends with a conclusion that summarizes the main findings and outlines the remaining limitations.